\documentclass[11pt]{article}
\usepackage[margin=1in]{geometry}
\usepackage{amsmath,amssymb,amsthm}
\usepackage{graphicx}
\usepackage[hidelinks]{hyperref}

\newtheorem{theorem}{Theorem}
\newtheorem{lemma}{Lemma}
\newtheorem{corollary}{Corollary}
\newtheorem{remark}{Remark}

\newcommand{\N}{\mathbb N}
\newcommand{\Kop}{\mathcal K}
\newcommand{\Lop}{\mathcal L}
\newcommand{\dens}{\operatorname{dens}}
\newcommand{\rc}{\operatorname{rc}}
\newcommand{\cardc}{\mathfrak c}
\newcommand{\norm}[1]{\left\lVert #1\right\rVert}

\title{An Exact Density Obstruction for Bounded Versus Precompact Sequences}
\author{FA Banach run \texttt{fa\_banach\_001}, agent \texttt{agent\_lane\_10}}
\date{23 July 2026}

\begin{document}
\maketitle

\section*{Status}

This packet records a candidate partial result, likely valid, toward Problem
5.1 and Remark 5.3 of Candido \cite{Candido}. It computes exactly the
densities of the two sequence spaces in the problem and derives a new
necessary cardinal condition for a positive example. It does not settle the
surviving cases.

\section*{Source Passage}

\begin{center}
\includegraphics[width=0.98\linewidth]{figures/open_problem_crop.png}
\end{center}

\section*{The Question and Its Intended Scope}

Candido asks whether there is a compact Hausdorff space \(K\) such that
\[
 C(K,\ell_\infty)\cong \ell_\infty(C(K)).
 \tag{1}\label{eq:source}
\]
Using \(\ell_\infty\cong C(\beta\N)\) and the exponential law for continuous
functions on compact spaces, the source identifies
\[
 C(K,\ell_\infty)
 \cong C(K\times\beta\N)
 \cong C(\beta\N,C(K)).
 \tag{2}\label{eq:exponential}
\]
Restriction from \(\beta\N\) to \(\N\) identifies
\(C(\beta\N,C(K))\) isometrically with the bounded sequences in \(C(K)\)
having relatively compact range. Thus the verbal formulation of Problem 5.1
asks whether all bounded \(C(K)\)-valued sequences, as a Banach space, can be
isomorphic to the precompact ones.

Problem 5.1 itself does not say ``infinite'', but the following paragraph,
Theorem 5.2, and Remark 5.3 explicitly ask for infinite \(K\). This restriction
is essential. If \(K\) is nonempty and finite, then \(C(K)\) is
finite-dimensional and both sides of \eqref{eq:source} are finite powers of
\(\ell_\infty\), hence are isomorphic. We address the substantive
infinite-\(K\) question.

\section*{Definitions and Proof Intuition}

Let \(X\) be a Banach space and let \(\Gamma\) be a set. Write
\[
 \ell_\infty(\Gamma,X)
 =
 \left\{x=(x_\gamma)_{\gamma\in\Gamma}:
       \sup_{\gamma\in\Gamma}\norm{x_\gamma}<\infty\right\}
\]
with the supremum norm, and put
\[
 \rc_\Gamma(X)
 =
 \left\{x\in\ell_\infty(\Gamma,X):
       \{x_\gamma:\gamma\in\Gamma\}\text{ is relatively compact in }X
 \right\}.
\]
The latter is closed and is therefore a Banach space. Equivalently,
\[
 \rc_\N(X)\cong C(\beta\N,X)\cong\Kop(\ell_1,X),
 \qquad
 \ell_\infty(\N,X)\cong\Lop(\ell_1,X).
 \tag{3}\label{eq:operators}
\]
The operator identifications follow by sending \(T\) to
\((Te_n)_{n\in\N}\).

The density of \(\rc_\Gamma(X)\) is controlled by finite-range maps, because
a relatively compact set has a finite net at every scale. By contrast,
\(\ell_\infty(\Gamma,X)\) remembers an independent choice at every
coordinate. A large separated subset of \(X\) therefore exponentiates to a
large separated subset of the full bounded-family space. These observations
give exact, rather than merely one-sided, density formulas.

\section*{Exact Density Formulas}

\begin{theorem}[Bounded versus precompact families]\label{thm:density}
Let \(X\) be an infinite-dimensional Banach space with
\(\dens X=\kappa\), and let \(\Gamma\) be an infinite set with
\(|\Gamma|=\lambda\). Then
\[
 \dens\rc_\Gamma(X)=\max\{\kappa,2^\lambda\}
 \tag{4}\label{eq:rc-density}
\]
and
\[
 \dens\ell_\infty(\Gamma,X)=\kappa^\lambda.
 \tag{5}\label{eq:bounded-density}
\]
\end{theorem}

\begin{proof}
Fix a norm-dense subset \(D\subset X\) of cardinality \(\kappa\).

We first prove \eqref{eq:rc-density}. Let
\(\mathcal F\) be the set of all finite-range maps from \(\Gamma\) to \(D\).
If \(x\in\rc_\Gamma(X)\) and \(\varepsilon>0\), compactness of the closure of
its range gives a finite \(\varepsilon/2\)-net. Approximate the points of that
net by points of \(D\) and assign each coordinate one of the resulting
approximants. This produces \(y\in\mathcal F\) with
\(\norm{x-y}_\infty<\varepsilon\). Hence \(\mathcal F\) is dense in
\(\rc_\Gamma(X)\).

There are \(\kappa\) finite subsets of \(D\). For each fixed nontrivial finite
range, the number of maps from \(\Gamma\) into that range is at most
\(2^\lambda\). Consequently
\[
 |\mathcal F|\le\max\{\kappa,2^\lambda\},
\]
which proves the upper bound in \eqref{eq:rc-density}.

The constant families form an isometric copy of \(X\) in
\(\rc_\Gamma(X)\), so its density is at least \(\kappa\). Fix a unit vector
\(x_0\in X\). The map
\[
 (a_\gamma)_{\gamma\in\Gamma}
 \longmapsto
 (a_\gamma x_0)_{\gamma\in\Gamma}
\]
is an isometric embedding of scalar \(\ell_\infty(\Gamma)\) into
\(\rc_\Gamma(X)\): a bounded subset of the scalar field has compact closure.
Finally,
\[
 \dens\ell_\infty(\Gamma)=2^\lambda.
\]
Indeed, the upper bound follows from cardinality, while
\(\{0,1\}^{\Gamma}\) is a \(1\)-separated family of cardinality
\(2^\lambda\). This proves the matching lower bound in
\eqref{eq:rc-density}.

For \eqref{eq:bounded-density}, coordinatewise approximation by \(D\) shows
that the bounded members of \(D^\Gamma\) form a dense subset of
\(\ell_\infty(\Gamma,X)\). Its cardinality is at most \(\kappa^\lambda\).

For the reverse inequality, we construct a \(1/2\)-separated subset
\(S\) of the unit sphere of \(X\) with \(|S|=\kappa\). By transfinite
recursion, suppose that \(x_\beta\) has been chosen for \(\beta<\alpha<\kappa\).
The closed linear span of the earlier vectors has density at most
\(\max\{|\alpha|,\aleph_0\}<\kappa\), so it is a proper closed subspace of
\(X\). Riesz's lemma gives a unit vector \(x_\alpha\) whose distance from
that subspace is greater than \(1/2\). The resulting family is
\(1/2\)-separated.

Every map from \(\Gamma\) into \(S\) is a bounded family, and two distinct
maps have supremum distance greater than \(1/2\). Thus
\(\ell_\infty(\Gamma,X)\) contains a uniformly separated set of cardinality
\[
 |S^\Gamma|=\kappa^\lambda.
\]
A metric space containing an \(\varepsilon\)-separated set of cardinality
\(\mu\) has density at least \(\mu\). This gives the lower bound and
completes the proof.
\end{proof}

\section*{Application to \(C(K)\)}

We use the standard equality
\[
 \dens C(K)=w(K)
 \tag{6}\label{eq:ck-density}
\]
for every infinite compact Hausdorff space \(K\). For completeness, here is
a proof. If \(D\) is norm dense in \(C(K)\), then \(D\) separates points:
approximate a continuous function separating any prescribed pair. The
evaluation map embeds \(K\) into the product of \(|D|\) copies of the scalar
field, so \(w(K)\le |D|\) for infinite cardinals. Conversely, a compact
Hausdorff space of weight \(w(K)\) has a point-separating family of continuous
scalar functions of that cardinality. The unital algebra over a countable
dense subfield generated by this family still has cardinality \(w(K)\), and
it is norm dense by the Stone--Weierstrass theorem \cite[Chapter 7]{Rudin}.
This proves \eqref{eq:ck-density}.

\begin{corollary}[Exact density test for the source problem]
\label{cor:ck}
Let \(K\) be an infinite compact Hausdorff space and put
\(\kappa=w(K)\). Then
\[
 \dens C(K,\ell_\infty)=\max\{\kappa,\cardc\},
 \qquad
 \dens\ell_\infty(C(K))=\kappa^{\aleph_0}.
 \tag{7}\label{eq:ck-formulas}
\]
In particular, the two spaces have the same density if and only if
\[
 \kappa^{\aleph_0}=\max\{\kappa,\cardc\}.
 \tag{8}\label{eq:density-equality}
\]
\end{corollary}

\begin{proof}
By \eqref{eq:exponential}, \(C(K,\ell_\infty)\) is the precompact-range
sequence space \(\rc_\N(C(K))\). Apply Theorem \ref{thm:density} with
\(\lambda=\aleph_0\), use \(2^{\aleph_0}=\cardc\), and invoke
\eqref{eq:ck-density}. The equality criterion is immediate.
\end{proof}

Candido's Theorem 5.2 supplies an additional necessary condition: if
\eqref{eq:source} holds for infinite \(K\), then \(\ell_\infty\) embeds
isomorphically into \(C(K)\). Since density is monotone for subspaces of
metric spaces,
\[
 \kappa=\dens C(K)\ge\dens\ell_\infty=\cardc.
\]
Combining this with Corollary \ref{cor:ck} gives the promised obstruction.

\begin{theorem}[Cardinal obstruction to a positive example]
\label{thm:obstruction}
If an infinite compact Hausdorff space \(K\) satisfies
\[
 C(K,\ell_\infty)\cong\ell_\infty(C(K)),
\]
then, with \(\kappa=w(K)\),
\[
 \boxed{\ \kappa\ge\cardc
 \quad\text{and}\quad
 \kappa^{\aleph_0}=\kappa.\ }
 \tag{9}\label{eq:boxed}
\]
Consequently no such \(K\) exists whenever
\(w(K)^{\aleph_0}>w(K)\). In particular, no example has weight of countable
cofinality.
\end{theorem}

\begin{proof}
The embedding conclusion of \cite[Theorem 5.2]{Candido} gives
\(\kappa\ge\cardc\). Isomorphic Banach spaces have the same density, so
\eqref{eq:ck-formulas} gives
\[
 \kappa^{\aleph_0}
 =
 \max\{\kappa,\cardc\}
 =
 \kappa.
\]

For the last assertion, suppose \(\operatorname{cf}(\kappa)=\aleph_0\).
Choose an increasing sequence of cardinals \(\kappa_n<\kappa\) cofinal in
\(\kappa\), and partition a set of cardinality \(\kappa\) into sets \(I_n\)
with \(|I_n|\le\kappa_n\). If the product
\(\prod_n\kappa_n^+\) had cardinality at most \(\kappa\), list all its
members as \((f_\xi)_{\xi<\kappa}\). At coordinate \(n\), fewer than
\(\kappa_n^+\) values occur among the \(f_\xi(n)\) with \(\xi\in I_n\);
choose \(g(n)\) outside those values. Then \(g\) differs from every
\(f_\xi\), a contradiction. Hence
\[
 \kappa
 <
 \prod_n\kappa_n^+
 \le
 \kappa^{\aleph_0},
\]
so \(\kappa^{\aleph_0}=\kappa\) is impossible.
\end{proof}

\section*{Limitations and Surviving Cases}

Theorem \ref{thm:obstruction} is a substantial necessary condition, not a
complete answer to Problem 5.1. It eliminates every weight on which
countable exponentiation strictly raises cardinality. It does not distinguish
the spaces when \(\kappa\ge\cardc\) and
\(\kappa^{\aleph_0}=\kappa\). In particular,
\[
 \cardc^{\aleph_0}
 =
 (2^{\aleph_0})^{\aleph_0}
 =
 2^{\aleph_0}
 =
 \cardc,
\]
so weight \(\cardc\) survives in ZFC.

At the operator level, the remaining question asks whether, for some
\(X=C(K)\) satisfying the stringent structural conditions in the source,
\[
 \Kop(\ell_1,X)\cong\Lop(\ell_1,X).
\]
The present density invariant cannot decide this fixed-point-cardinal regime.
No positive example or full nonexistence proof is claimed.

\section*{Novelty Check}

The run indexes were searched for the arXiv id, the exact paper title,
bounded versus precompact \(C(K)\)-valued sequences, both sides of
\eqref{eq:source}, and the operator formulation \eqref{eq:operators}. No
overlapping packet was found.

Bounded external searches on 23 July 2026 used the exact phrases
\begin{quote}
\texttt{"space of bounded sequences in C(K)" "precompact sequences" isomorphic},\\
\texttt{"C(K,ell\_infty)" "ell\_infty(C(K))"},\\
\texttt{"K(ell\_1,X)" "L(ell\_1,X)" isomorphic},
\end{quote}
as well as the paper title and citation searches. They found the arXiv and
published versions of \cite{Candido}, but no later paper explicitly answering
the problem or recording the density obstruction. This bounded search is
evidence, not a guarantee, of novelty.

\section*{Human-Review Recommendation}

The proof is elementary and self-contained after the source's Theorem 5.2 and
the standard Stone--Weierstrass theorem. The packet should be reviewed as an
exact density characterization and a new necessary condition. It should not
be promoted to a full solution unless the surviving
\(\kappa^{\aleph_0}=\kappa\) cases are resolved by an additional Banach-space
invariant.

\begin{thebibliography}{9}

\bibitem{Candido}
L. Candido,
\emph{Complementations in \(C(K,X)\) and \(\ell_\infty(X)\)},
Colloquium Mathematicum \textbf{172} (2023), 129--141;
arXiv:2104.07152.
\href{https://doi.org/10.4064/cm8868-10-2022}
{doi:10.4064/cm8868-10-2022}.

\bibitem{Rudin}
W. Rudin,
\emph{Real and Complex Analysis}, third edition,
McGraw--Hill, 1987.

\end{thebibliography}

\end{document}

